%% ---------------------------------------------------------------------
\section{Introduction}\label{sec:intro}

A bank of Gabor filters, fixed before training and free to compute, is
worth $+26$ accuracy points to a ViT-B/16 at $224$\,px. Injected exactly
the same way into a network that has already had contrastive
pre-training, it is worth nothing. Injected into one initialized from
ImageNet, it costs $17$ points. One intervention, one frozen recipe, one
code path: three outcomes spanning forty accuracy points. The useful
question is not which number is representative, because none of them is.
It is what distinguishes the three cases, and whether that can be
determined \emph{before} the combination is trained.

This is the general situation, not a curiosity of one prior. Every
practitioner who has trained a network on a small dataset faces the same
menu of remedies: pre-train on ImageNet, run a self-supervised stage,
augment more aggressively, or inject prior knowledge into the model. Each
fuses an external source of information (weights learned elsewhere,
invariances distilled from augmented views, hand-crafted structure) with
whatever the labeled data can teach on its own. What the field lacks is
not remedies but a controlled account of how these sources
\emph{combine}: when does adding a second source help, when is it
redundant with what the first already supplies, and when does the fusion
actively destroy value? These three outcomes, which we call
\emph{stack}, \emph{substitute} and \emph{tax}, are usually discovered
anecdotally, one paper and one setting at a time, under training recipes that
differ in a dozen uncontrolled ways.

Our answer is that sources trade in identifiable \emph{currencies}, and
that what a source is worth depends on which currency is scarce rather
than on how modern the method is. Two sources add only when their
currencies differ. This is not a metaphor we impose after the fact: each
source's currency is measurable, on frozen features, from each source
\emph{alone}, before any combination is trained. That is what makes the
three outcomes above predictable instead of anecdotal, and it is why the
same free prior can be the most valuable intervention we measure in one
setting and the most damaging in another.

This paper treats the question as a measurement problem. We build a single
instrument (one frozen training recipe, fixed data subsets shared by every
run, and a numerically pinned bank of Gabor and moment spectral filters) and
use it to measure one deliberately austere source of prior knowledge against
representative data-driven alternatives at matched or declared compute. The prior, which we call \emph{MomentAux}, could hardly cost
less: a fixed bank of oriented-energy targets, injected only during training
through a small auxiliary regression head whose loss weight decays to exactly
zero, costing ${\sim}2\%$ extra training compute and \emph{nothing} at
inference. Against it we field SimCLR, SimSiam, and DINO initializations
(each ${\sim}2\times$ compute), ImageNet transfer, DeiT-strength
augmentation, learned FitNets-style teachers, and targets built from histograms of oriented gradients (HOG). The
resulting grid spans 13 datasets across six visual domains, nine backbone
families from ResNet-18 to ViT-B/16, data scales from 150 to 1.28M images,
resolutions from $32$ to $224$\,px, and \computeCells{} configurations trained
over \computeRuns{} runs, each with a retained run record.

A grid of this size threatens to be a table dump. What rescues it is an
empirical law we found rather than assumed, then tested by registering its
predictions in advance and trying to break them. Writing $\Delta$ for an intervention's end-to-end
accuracy gain and $G$ for the gain a linear evaluation reads from its frozen
features under identical probing, we find across the grid that
\begin{equation}
  \Delta \;=\; G \;+\; \mathrm{readout}(\mathrm{base}),
  \label{eq:law}
\end{equation}
where the readout term depends, to a good approximation, only on the
\emph{baseline accuracy} of the cell: it is negative below a crossing
located at ${\sim}32$--$40\%$ accuracy, near zero at the crossing, and
small and positive above it. The sign prediction holds in $\auditRate\%$ of the
\auditResolvable{} cells where measurement uncertainty allows it to be tested, and
in $94\%$ of those below the crossing
(Section~\ref{sec:law}), and its quantitative form survived our hardest
pre-registered tests: it predicted the feature gain of a backbone family it
had never seen (Swin) from baseline heights alone, and at ImageNet scale it
tied end-to-end gains to feature gains within $1.1$ accuracy points on five
of six backbone pairs (Section~\ref{sec:scale}).

The law converts the grid into an explanation, and it is where currency stops
being a metaphor and becomes the term $G$: heavy augmentation supplies
nuisance-invariance, contrastive pre-training much the same invariance at
greater expense, ImageNet weights mature photographic features, and the moment
prior oriented-energy structure.
Same-currency fusions substitute; the prior and SimCLR each render the
other nearly worthless, and their combination never beats the better
single. Different-currency fusions stack; the prior's gain under
DeiT-strength augmentation does not shrink but grows by a factor of
$1.4$--$2.4$. And fusing a shaping prior into an initialization that
already carries mature features taxes accuracy in proportion to what the
initialization supplied, a cost we show lands almost entirely on the features themselves, not the classifier
(Section~\ref{sec:fusion}).

Two headline findings illustrate what the instrument can resolve. First,
small vision transformers carry a large feature deficit that the prior
fills: the gain reaches $+13.0$ points for ViT-S/16 and $+26.0$ for
ViT-B/16 on ImageNet-100 at $224$\,px under the standard DeiT recipe,
\emph{growing} with model scale, and persists at $+3.2$ after 1.28 million
images where the ResNet baseline is exactly neutral ($+0.04$). Second, the fashionable remedy is not always the
right one: under plain augmentation SimCLR beats the prior across the
mid-data band on photographic datasets, but under the DeiT recipe the
comparison flips at every fraction at matched ${\sim}2\times$ cost, because the
recipe supplies the invariance SimCLR was paid that compute to learn. At
$5\times$ the flip narrows to the higher-data band, where the prior still wins
on both populations, while the comparator takes the scarcest cells
(Section~\ref{sec:budget}). The ordering tracks the data fraction, so this is
a statement about budget and regime, not a ranking of methods.

\paragraph{Contributions.}
\begin{enumerate}
  \item \textbf{An instrument, not just experiments} (Section~\ref{sec:instrument}):
    a frozen recipe with fixed subsets, pinned filter banks, declared
    cost normalization, and pre-registered predictions with falsifiers,
    every headline table regenerable by one command from released artifacts.
  \item \textbf{The grid} (Section~\ref{sec:grid}): $\computeCells$ configurations,
  $\computeRuns$ runs,
    comparing a free spectral prior against SSL, transfer, augmentation and
    learned teachers across 13 datasets, 9 backbones and
    150 to 1.28M images.
  \item \textbf{The law} (Section~\ref{sec:law}): $\Delta = G +
    \mathrm{readout}(\mathrm{base})$, validated from CIFAR scale to
    ImageNet scale, with its boundary cases reported as findings and not excluded, including the
    one patterned failure where features improve at sufficiency while accuracy
    does not follow (Section~\ref{sec:exceptions}). It also resolves an
    apparent puzzle at ImageNet scale, where three backbones on identical data
    give three different envelope shapes (Section~\ref{sec:scaleenvelope}).
  \item \textbf{A fusion taxonomy with mechanism}
    (Section~\ref{sec:fusion}): stack, substitute, or tax, decided by
    whether the fused sources carry the same information currency, and
    verified on the feature side, not just end-to-end.
  \item \textbf{A procedure, not only a taxonomy}
    (Section~\ref{sec:procedure}): because each source's currency is
    measurable from that source \emph{alone}, the outcome of combining two
    is predictable before the combination is trained. We state this as an
    algorithm, run it end to end on a pair whose answer we then verify by
    training, and distil the grid into a regime-indexed guide with costs
    attached (Section~\ref{sec:guide}).
\end{enumerate}

In one line, what the grid says: measure each source's currency alone and the
outcome of fusing them can be anticipated in advance
(Section~\ref{sec:fusion}); do not pour a shaping prior into a strong
initialization, because under the schedule we test the tax scales with what
that initialization was worth (Section~\ref{sec:taxsub}); modern augmentation does not replace
structure but amplifies it (Section~\ref{sec:stack}); attention's feature
deficit grows with model scale rather than shrinking
(Section~\ref{sec:scale}); and a comparator's budget is part of the claim,
since two of ours did not survive raising it
(Section~\ref{sec:budget}).

%% ---------------------------------------------------------------------
\section{Related work}\label{sec:related}

\looseness=-1 We group prior work by the mechanism through which it supplies
information to a data-limited network, since that is the axis along which
our comparison is organized: hand-crafted structure placed in the forward
path; auxiliary and distillation targets that shape features during
training; self-supervised and transfer initializations that pay compute
for representation; augmentation recipes that supply invariance;
small-data results specific to attention; and the large-scale comparative
benchmarks that most closely resemble what we build. Each paragraph below
takes one of these in turn. Section~\ref{sec:fusiontheory} then places our
three outcomes against the fusion literature's own account of how sources
relate, and Section~\ref{sec:gap} states, in a single table, which
properties a controlled fusion study requires and which of them each line
of work currently provides.

\paragraph{Hand-crafted structure inside learned networks.}
Fixed oriented filters predate deep learning as image descriptors
\citep{dalal2005hog}. Scattering networks showed that
cascades of fixed wavelets rival learned features in small-data regimes
\citep{bruna2013scattering}, and Gabor-structured
convolutions have been used to regularize or cheapen early layers
\citep{luan2018gaborcnn}. These lines place the structure
\emph{in the forward path}, where it occupies input bandwidth permanently.
Our forward-path baselines reproduce their known failure mode, a
mid-data penalty band that no filter choice escapes
(Section~\ref{sec:ablations}), and our central object differs precisely
in that the structure lives in a training-only auxiliary target whose
influence decays to zero, so abundant data can override the prior rather
than pay for it.

\paragraph{Auxiliary targets, deep supervision, and distillation.}
\looseness=-1 Regressing intermediate features onto targets is the mechanism of FitNets
\citep{romero2015fitnets} and deep supervision \citep{lee2015deeply}; HOG
targets power MaskFeat's self-supervised pre-training
\citep{wei2022maskfeat}. What distinguishes our setting is the
\emph{fixedness and cost} of the target: no teacher network is trained, no
pre-training stage is run, and the target is a closed-form function of the
input. Our controls quantify how much this matters: a learned same-data
teacher moves accuracy by at most $0.7$ points anywhere on an eight-point
envelope, HOG targets recover a third to a half of the moment target's gain, and a
random fixed target does nothing (Table~\ref{tab:controls}), so the gain
is attributable to the moment structure itself, not to auxiliary
regression as a mechanism.

\paragraph{Self-supervised learning (SSL) as initialization.}
Contrastive and self-distilling pre-training
\citep{chen2020simclr,chen2021simsiam,caron2021dino}
is the modern default source of ``free'' features, and its large-scale
recipes are well established. Much less is measured about the small-data,
from-scratch regime we target. Our
grid contributes three facts to that gap: negative-free SimSiam learns
essentially nothing at $2.5$--$25$k images under its published CIFAR
recipe, though four times that budget recovers up to $+4.3$ and we report
it as a property of the budget rather than of the method; SimCLR is strong
on photographic statistics but loses to a free domain-agnostic prior on
satellite imagery in the low-data band; and the SSL-versus-prior comparison
inverts under a modern supervised recipe at matched cost, because the
augmentation stack substitutes for the invariance SSL sells
(Sections~\ref{sec:fusion} and \ref{sec:budget}).

\paragraph{Transfer learning.}
ImageNet initialization remains the strongest single intervention on
photo-like data, with known limitations under
domain shift and with training-from-scratch competitive given enough data
\citep{he2019rethinking}. We use transfer both as a comparator and as a
test of fusion: injecting the spectral prior \emph{on top of} ImageNet
weights is the one configuration in our grid that is destructive
everywhere, and linear probing shows the destruction is feature-side,
early shaping erases exactly the mature structure the initialization
supplied, in proportion to how much that structure was worth
(Section~\ref{sec:fusion}).

\paragraph{Small-data vision transformers.}
That ViTs underperform without large data or heavy regularization is well
documented \citep{dosovitskiy2021vit,touvron2021deit},
and a family of remedies exists: distillation tokens
\citep{touvron2021deit}; architectures that reintroduce the hierarchical,
local inductive bias attention lacks, such as Swin \citep{liu2021swin},
which we include as a backbone for that reason; and auxiliary
self-supervised objectives trained jointly with the classifier, of which
\citet{liu2021efficientvit} is the closest antecedent to this work: a
dense relative-localization task added for exactly this deficit, differing
from ours in that its target is computed from the images themselves rather
than pinned in advance. Our contribution
to this literature is quantitative and comparative: we measure the deficit
as a feature quantity ($G \approx 13$--$15$ points for ViT-tiny across
$2.5$--$12.5$k images, roughly $2.3\times$ any convolutional backbone),
show a fixed spectral target fills it at lower cost than contrastive
pre-training and more effectively than DeiT augmentation alone, and show
the deficit \emph{grows} with model scale at fixed data
(Section~\ref{sec:scale}), which argues it reflects data-hunger, not small-model capacity.

\paragraph{Large-scale comparative benchmarks.}
Two recent efforts share our comparative ambition and clarify what is
still missing. \emph{Battle of the Backbones}
\citep{goldblum2023battle} compares a broad suite of pretrained
backbones across tasks over more than 1{,}500 training runs, and finds
supervised convolutional pre-training still competitive with
self-supervised and vision-language alternatives. \citet{marks2025closer}
take a closer look at self-supervised benchmarking specifically, and show
that conclusions depend heavily on the evaluation protocol chosen. Both
benchmark \emph{pretrained backbones}, so their unit of comparison is a
checkpoint someone else produced, and neither normalizes by training
cost nor asks what happens when two sources of information are combined.
Our unit is instead the \emph{intervention applied to a fixed recipe},
which is what makes cost normalization and the stacking question
answerable.

\paragraph{Spectral structure inside transformers.}
A parallel line injects spectral structure into attention architectures:
FViT \citep{shi2026fvit} adds a \emph{learnable} Gabor filter to focus
attention across scales and orientations, and scattering-based token mixers
do the same with wavelet decompositions. These are architectural changes that
persist at inference, and each is evaluated as a model proposal against
published baselines. Our object is the
complement: the filters are fixed rather than learned, they never enter
the forward path, and the deployed network is unchanged, which is what
lets us ask the benchmark question of when adding them pays rather than
the model question of whether one architecture beats another.


\subsection{Relation to fusion theory}\label{sec:fusiontheory}
The outcomes we measure are not new categories, and we set out the prior
claims before our own. \citet{durrantwhyte1988sensor} classified the
interactions between disparate information sources as \emph{cooperative},
\emph{competitive} and \emph{complementary}; the surveyed form now standard
in the field \citep{khaleghi2013multisensor} states the triple as
complementary, redundant and cooperative, where complementary sources
describe different aspects of the target, redundant sources describe the
same aspect, and cooperative sources jointly yield what neither yields
alone. Our \emph{stack} is complementary combination and our
\emph{substitute} is redundant combination. We adopt that vocabulary here
and use our own terms only as informal labels for it.

\citet{dasarathy2001what} organizes the field's questions as what, where,
why, when and how to fuse, and observes that the \emph{when} has received
least attention because it is the one that cannot be answered from the
architecture. That is the question this paper takes up, in the specific
setting where one source is hand-crafted and the other is learned from the
data. Nor is knowledge fusion itself a new object: \citet{smirnov2019knowledge}
surveys the patterns by which knowledge and data are combined, and catalogs
the same three outcomes we recover. What that survey records as patterns to
be chosen by a designer, we measure as outcomes to be predicted from the
sources, which is the difference between a taxonomy and a decision
procedure. The field's benchmarks have generally compared \emph{algorithms}
for a fixed fusion problem; ours compares
\emph{sources} under a fixed algorithm, holding the recipe frozen so that
what varies is the knowledge injected, not the method injecting it.

One difference is substantive, not terminological, and it sharpens
what a representation-learning setting adds. In sensor fusion redundancy is
a \emph{virtue}: two instruments measuring the same quantity reduce
variance and buy fault tolerance. In ours it is waste, because performance
is bounded by the information the sources share rather than improved by
supplying it twice. Redundancy pays for reliability and not for
representation quality, and a taxonomy built for the first case does not
transfer its sign to the second.

Our third outcome, in which a source destroys what its partner supplied,
is likewise not new, and has at least three names. It is \emph{catastrophic
fusion} in the sense of \citet{movellan1997catastrophic}, where modules
combined under mismatched assumptions do worse than the better module
alone; it is a form of negative transfer
\citep{wang2019negative}; and on the multimodal side it
is the modality competition of \citet{wang2020multimodal}, who report that
the best unimodal network often beats the jointly trained multimodal one
despite the latter receiving strictly more information. What we add is not
the phenomenon but its magnitude law: the damage scales with how much the
displaced source had contributed, and vanishes exactly where it had
contributed nothing.

The formal account of our informal ``currency'' is partial information
decomposition \citep{williams2010nonnegative}, which splits what two
sources tell us about a target into unique, redundant and synergistic
atoms, and which \citet{liang2023quantifying} have already carried into
deep multimodal learning with estimators that scale and with a-priori model
selection as the application. We do not compute it. Our measurement is a
linear-evaluation gap on frozen features: inexpensive enough for \computeCells{} cells, ordinal, not an information quantity, and blind to anything not
linearly decodable. Read our result as an empirical instance of the
decomposition those authors formalize, not as an estimate of it.

Two further connections position the predictor, not the taxonomy.
Predicting a downstream outcome from an inexpensive measurement, without training
the combined system, is the programme of transferability estimation
\citep{nguyen2020leep}, which scores a single pre-trained
source; ours asks the same question of a \emph{pair}. And the challenge
that programme raises for us is direct:
\citet{newell2020useful} report that linear evaluation does not correlate
with fine-tuning performance in general. Our defence is narrow and empirical, not a rebuttal. We do not use the linear evaluation to
rank absolute quality; we use the \emph{gap} between two arms probed under
an identical protocol, on checkpoints that differ in one intervention, and
we report how well that gap predicts the paired difference
(Section~\ref{sec:law}) rather than assuming it does.

Finally, our knowledge source sits where \citet{vonrueden2023informed} make
precise. Their taxonomy classifies informed learning by knowledge
\emph{source}, \emph{representation} and \emph{integration point}, and
states that informed learning requires two independent sources of
information, data and prior knowledge, which is the reading of ``two
sources'' under which this study is a fusion study. Ours is
scientific knowledge about early vision, represented as a pinned filter
bank, integrated at the learning algorithm; the nearest existing method is
their own informed pre-training, which
condenses knowledge into prototypes for initialization rather than into a
decaying auxiliary target. That survey catalogs where knowledge can be
injected. What it does not settle, and what a practitioner must decide, is
\emph{when} injecting it is worth anything given what the data and the
available pre-trained artifacts already supply. That is the question this
benchmark is built to answer, and we are careful to claim only that: we
found no controlled, compute-declared answer to it, not that no neighbouring
question has been answered.

\subsection{The gap this paper closes}\label{sec:gap}

\begin{table}[pos=htbp]
\caption{Where the literature stands on the five properties a fusion
question needs, and what this work adds. \cmark{} = property present;
\pmark{} = partial or incidental; \xmark{} = absent. ``Cost-normalized''
means comparators are reported against a declared training-compute
multiple rather than at whatever budget each method's authors chose;
``feature-side'' means the effect is measured on frozen representations,
not only end-to-end; ``combination'' means the study measures what
happens when two sources are applied together.}
\label{tab:gap}
\centering\scriptsize
\setlength{\tabcolsep}{2.6pt}
\zebra{2}
\begin{tabular}{lccccc}
\toprule
Line of work & \rotbox{one frozen recipe} & \rotbox{cost-normalized} &
  \rotbox{feature-side} & \rotbox{combinations} & \rotbox{predictive model} \\
\midrule
Fixed spectral filters in the forward path
  \citep{bruna2013scattering,luan2018gaborcnn,shi2026fvit}
  & \pmark & \xmark & \xmark & \xmark & \xmark \\
Auxiliary and distillation targets
  \citep{romero2015fitnets,lee2015deeply,wei2022maskfeat}
  & \pmark & \xmark & \pmark & \xmark & \xmark \\
Self-supervised pre-training
  \citep{chen2020simclr,chen2021simsiam,caron2021dino}
  & \xmark & \xmark & \cmark & \xmark & \xmark \\
Transfer learning studies
  \citep{he2019rethinking}
  & \pmark & \xmark & \pmark & \xmark & \xmark \\
Small-data ViT recipes
  \citep{touvron2021deit}
  & \xmark & \pmark & \xmark & \pmark & \xmark \\
Backbone and self-supervised benchmarks
  \citep{goldblum2023battle,marks2025closer}
  & \pmark & \xmark & \cmark & \xmark & \xmark \\
Classical fusion taxonomy
  \citep{durrantwhyte1988sensor,williams2010nonnegative}
  & \xmark & \xmark & \xmark & \cmark & \pmark \\
\midrule
\textbf{This work} & \cmark & \cmark & \cmark & \cmark & \cmark \\
\bottomrule
\end{tabular}
\end{table}

Table~\ref{tab:gap} summarizes the position. Each column is a property
that a question about \emph{fusing} information sources requires, and no
existing line supplies all five at once.

The first column matters because comparisons made under each method's own
preferred recipe cannot attribute a difference to the method. Benchmarks
of pretrained backbones \citep{goldblum2023battle,marks2025closer} inherit
whatever recipe produced each checkpoint, and \citet{marks2025closer}
show directly that the evaluation protocol can change the conclusion. We
instead hold one recipe and one set of fixed image indices fixed and
vary only the intervention.

The second is rarely reported. Self-supervised pre-training costs
roughly twice the training compute of the baseline it improves, but that
factor almost never appears next to the accuracy it buys, which makes
``method A beats method B'' unanswerable for a practitioner with a fixed
budget. Every comparator here carries a declared multiple.

The third and fourth columns are what turn a table of numbers into an
explanation. Linear probing is standard practice for evaluating
self-supervised representations, so the feature-side view exists in that
literature; what is missing is applying it to a \emph{paired} comparison
so the same intervention is measured end-to-end and on frozen features at
once, which is what exposes redundancy between two sources. Combinations
are the least studied of all: the question of what happens when a prior
is applied on top of augmentation, on top of self-supervised weights, or
on top of a pretrained initialization is precisely the fusion question,
and we could find no controlled study that measures all three.

The fifth column is the one we regard as the paper's distinguishing
contribution. Benchmarks describe; a law predicts. Our decomposition
$\Delta = G + \mathrm{readout}(\mathrm{base})$ was registered in advance and
used to call all three classes of unmeasured quantity listed above before
those measurements existed (Sections~\ref{sec:law} and~\ref{sec:scale}). A related
decomposition of learning dynamics into representation and decoding components
has been proposed for grokking \citep{liu2022grokking}; to our knowledge no
comparable predictive account exists for data-efficiency interventions.

\paragraph{Benchmark methodology.}
Uncontrolled recipes manufacture illusory progress, and reproducibility in
machine learning remains limited. We adopt
the strictest discipline we could operationalize: a single frozen recipe;
fixed per-subset indices so every intervention sees byte-identical
images; numerically fingerprinted filter banks; deviations quarantined
marked as such and never mixed into headline tables; and,
unusual in this literature, predictions with explicit falsifiers
registered \emph{before} each wave of runs, with misses reported as
plainly as hits (Sections~\ref{sec:instrument}, \ref{sec:scale}).
